\chapter{Implementation Details}
This section describes the practical implementation of the proposed framework. I first detail the baseline model, including its architecture, training procedure, and 
inference pipeline. All implementation choices are reported to ensure reproducibility of the experiments. All code is implemented in Python using the PyTorch library.
All the hyperparameters are selected based on validation performance using the library Optuna and kept constant across all experiments unless otherwise specified. 
The experiments were conducted on a machine with TODO

\section{Baseline Implementation}
The \hyperref[sec:baseline]{baseline model} serves as the reference architecture upon which the proposed method is built. It takes as input a temporal sequence of 2D human poses $[B, T, J, 3]$,
where $B$ is the batch size equal to 32, $T$ is the number of frames in the sequence equal, and $J$ is the number of joints equal to 17, and 3 represents the $(x, y, confidence)$ coordinates of each 
joint and predicts the corresponding 3D pose representation, giving as output a tensor of shape $[B, T, J, 3]$, where the last dimension represents the $(x, y, z)$ 
coordinates of each joint in 3D space.

\subsection{Architecture}
The baseline model is based on a Transformer-based architecture designed to model temporal dependencies across pose sequences. Each input pose is first embedded into a 
latent feature space of dimension 512. The resulting embeddings are then processed by a stack of Transformer encoder layers.

Specifically, the Transformer module consists of 3 encoder layers, each employing multi-head self-attention with 8 attention heads. This design allows the model to 
capture both short-term and long-range temporal correlations within the input sequence. The output of the Transformer is finally mapped to the target pose 
representation through a linear projection layer.

No architectural modifications are introduced with respect to the baseline configuration, which is kept fixed throughout all experiments to provide a fair comparison 
with the proposed approach.

\begin{table}[t]
    \centering
    \caption{Baseline model architecture. The network takes as input a temporal sequence of 2D poses and outputs the corresponding pose representation.}
    \label{tab:baseline_architecture}
    \renewcommand{\arraystretch}{1.2}
    \begin{tabular}{lccc}
    \hline
    \textbf{Layer} & \textbf{Input Shape} & \textbf{Output Shape} & \textbf{Params} \\
                   & $(32, 243, 17, 3)$ & $(32, 243, 17, 3)$ &  \\
    \hline
    Linear (Input Embedding) & $(32, 243, 51)$ & $(32, 243, 512)$ & 26,624 \\
    \hline
    Positional Encoding & $(32, 243, 512)$ & $(32, 243, 512)$ & -- \\
    Transformer Encoder (3 layers, 8 heads) & $(32, 243, 512)$ & $(32, 243, 512)$ & 9,457,152 \\
    \hline
    Linear & $(32, 243, 512)$ & $(32, 243, 512)$ & 262,656 \\
    Tanh & $(32, 243, 512)$ & $(32, 243, 512)$ & -- \\
    Linear (Pose Head) & $(32, 243, 512)$ & $(32, 243, 51)$ & 26,163 \\
    \hline
    \textbf{Total} &  &  & \textbf{9,772,595} \\
    \hline
    \end{tabular}
\end{table}


\subsection{Training Procedure}
The baseline model is trained using the MPJPE as the optimization objective. Training is initially performed on the SportsPose dataset, which provides labeled pose 
sequences suitable for learning temporal motion patterns. Subsequently, the model is fine-tuned on the AirPose dataset to adapt it to the target domain.

The training process is carried out for a maximum of 100 epochs. The Adam optimizer is employed with an initial learning rate of $5 \times 10^{-4}$ and a weight decay of 
0.01. The learning rate is scheduled using a exponential decay with a factor of 0.99 applied every epoch.
A warm-up phase of 2 epochs is applied at the beginning of training to stabilize optimization. Early stopping is used with a patience of 9 epochs based on 
validation performance.

\subsection{Inference Pipeline}
During inference, the input to the system is a RGB video sequence represented as a tensor of shape $[T, H, W, 3]$, where 
For each frame, a human detection stage is first performed using YOLO to localize the subject of interest.

The detected human bounding box is cropped and resized to a fixed resolution of $256 \times 192$ pixels. The resulting image is then passed to a ViTPose model to 
perform 2D pose estimation, producing a set of 2D keypoints for each frame.

The extracted 2D keypoints are normalized to the range $[-1,1]$ using the original frame dimensions.
These normalized 2D poses are finally provided as input to the baseline network, which predicts the corresponding output pose sequence.
